% ============================================================
%  cap05/5.5_discusion_presentacion.tex
% ============================================================
\section{Discusi\'on de Resultados}

La presente investigaci\'on ha desarrollado, evaluado y desplegado experimentalmente el proyecto \textit{Democra}, una plataforma SaaS dise\~nada espec\'ificamente para absorber la carga operativa y documental de las Organizaciones No Gubernamentales (ONGs) en un entorno digital centralizado, concurrente y criptogr\'aficamente aislado. Los resultados emp\'iricos obtenidos en el Cap\'itulo 5 (Pruebas de Estr\'es, Evaluaci\'on de Sobrecarga de \textit{Row-Level Security}, Auditor\'ia de Intrusiones Trans-Tenant, cobertura de c\'odigo, pruebas unitarias y documentaci\'on de la API) proveen un espectro cuantitativo robusto que confirma la viabilidad arquitect\'onica del paradigma propuesto. En este cap\'itulo, dichos hallazgos se deconstruyen y se contrastan dial\'ecticamente con la literatura acad\'emica preexistente y los antecedentes expuestos en el Marco Te\'orico (Cap\'itulo 2), con el prop\'osito de extraer deducciones inferenciales sobre el estado del arte de la ingenier\'ia web aplicada al sector social no lucrativo.
